% Part: first-order-logic
% Chapter: syntax-and-semantics
% Section: Structures

\documentclass[../../../include/open-logic-section]{subfiles}

\begin{document}

\olfileid{fol}{syn}{str}

\olsection{د لومړۍ درجې ژبو \printtoken{P}{structure}}

\begin{explain}
د لومړۍ درجې ژبې پخپله \emph{بې‌تفسيره} دي: له !!{constant}s،
!!{function}s او !!{predicate}s سره کومه ځانګړې معنا نۀ وي تړلې۔
معنا د \article{structure} \emph{!!{structure}} په ټاکلو ورکول کېږي۔
هغه \emph{د تعريف ساحه} ټاکي؛ يعنې هغه څيزونه چې !!{constant}s ئې
ټاکي، !!{function}s پرې عمل کوي، او سورونه پرې ګرځي۔ سربېره پر دې،
دا هم ټاکي چې کوم !!{constant}s کوم څيزونه ټاکي، !!a{function}
څيزونه څنګه نورو څيزونو ته لېږي، او !!{predicate}s پر کومو څيزونو
صدق کوي۔ !!^{structure}s په منطق کښې د \emph{معنايي} مفهومونو بنسټ
دي، لکه لازميت، اعتبار، او د صدق وړتيا۔ په ليکنو کښې ورته په بېلابېل
ډول «جوړښتونه»، «تفسيرونه» يا «مدلونه» ويل کېږي۔
\end{explain}

\begin{defn}[!!^{structure}s]
د لومړۍ درجې منطق د يوې ژبې $\Lang{L}$ دپاره \Article{structure}
\emph{!!{structure}}~$\Struct M$ له لاندې اجزاوو جوړ دے:
\begin{enumerate}
\item \emph{د تعريف ساحه:} يو ناتش سټ، $\Domain M$
\item \emph{د !!{constant}s تفسير:} د $\Lang{L}$ د هر
  !!{constant}~$c$ دپاره يو !!a{element} $\Assign{c}{M} \in \Domain M$
\item \emph{د !!{predicate}s تفسير:} د $\Lang{L}$ د هر $n$-ځايه
  !!{predicate}~$R$ دپاره (له $\eq$ پرته) يوه $n$-ځايه اړيکه
  $\Assign{R}{M} \subseteq \Domain{M}^n$
\item \emph{د !!{function}s تفسير:} د $\Lang{L}$ د هر $n$-ځايه
  !!{function}~$f$ دپاره يوه $n$-ځايه تابع $\Assign{f}{M}
  \colon \Domain{M}^n \to \Domain{M}$
\end{enumerate}
\end{defn}

\begin{ex}
د حساب د ژبې دپاره !!^a{structure}~$\Struct M$ له يوه سټ؛ د
$\Domain M$ له يوه غړي $\Assign{\Obj 0}{M}$، چې د
!!{constant}~$\Obj 0$ تفسير دے؛ له يوې يو ځايه تابع
$\Assign{\Obj \prime}{M} \colon \Domain{M} \to \Domain M$؛ له دوو
دوه ځايه تابعو $\Assign{\Obj +}{M}$ او $\Assign{\Obj
  \times}{M}$، چې دواړه $\Domain M^2 \to \Domain M$ دي؛ او له يوې
دوه ځايه اړيکې $\Assign{\Obj <}{M} \subseteq \Domain{M}^2$ جوړ دے۔

د داسې جوړښت يوه څرګنده بېلګه دا ده:
\begin{enumerate}
\item $\Domain N = \Nat$
\item $\Assign{\Obj 0}{N} = 0$
\item د هر $n \in \Nat$ دپاره $\Assign{\Obj \prime}{N}(n) = n + 1$
\item د هر $n, m \in \Nat$ دپاره $\Assign{\Obj +}{N}(n, m) = n + m$
\item د هر $n, m \in \Nat$ دپاره $\Assign{\Obj \times}{N}(n, m) = n\cdot m$
\item $\Assign{\Obj <}{N} = \Setabs{\tuple{n, m}}{n \in \Nat, m \in
  \Nat, n < m}$
\end{enumerate}
د $\Lang L_A$ دپاره په دې ډول تعريف شوے جوړښت~$\Struct N$ د
\emph{حساب معياري مدل} بلل کېږي، ځکه د~$\Lang L_A$ غير منطقي ثابتونه
کټ مټ هغسې تفسيروي چې تمه ئې کوئ۔

خو د~$\Lang L_A$ دپاره ډېر نور ممکن !!{structure}s هم شته۔ مثلاً د
ساحې په توګه د~$\Nat$ پر ځاے د صحيح عددونو سټ~$\Int$ اخيستلے، او د
$\Obj 0$، $\Obj \prime$، $\Obj +$، $\Obj \times$ او $\Obj <$ تفسيرونه
ورسره سم تعريفولے شو۔ خو د~$\Lang L_A$ داسې جوړښتونه هم تعريفولے شو
چې له عددونو سره هېڅ لرې اړيکه هم نۀ لري۔
\end{ex}

\begin{ex}
د سټ تيورۍ د ژبې~$\Lang L_Z$ يو جوړښت~$\Struct M$ يوازې يوه ساحه
او يوه دوه ځايه اړيکه غواړي۔ نو په فني لحاظ مثلاً د خلکو سټ او اړيکه
«$x$ له $y$ څخه مشر دے» هم د $\Lang L_Z$ دپاره !!a{structure}
کېدے شي، او هم~$\Nat$ له $n \ge m$ سره د $n, m \in \Nat$ دپاره۔
\textbf{د ژباړې د نسخې سمون (OLFOL-012):} سرچينه د اړيکې د ځاييز
والي په عبارت کښې ناسمه نښلوونکې کرښه ږدي؛ دلته مطلب «يوه دوه ځايه
اړيکه» لوستل شوے دے۔

د $\Lang L_Z$ يو په ځانګړي ډول په زړه پورې !!{structure}، چې د ساحې
!!{element}s ئې په رښتيا سټونه او د $\Obj \in$ تفسير ئې په رښتيا
اړيکه «$x$ د~$y$ !!a{element} دے» وي، د \emph{موروثي متناهي
سټونو} !!{structure}~$\Struct{{HF}}$ دے:
\begin{enumerate}
\item $\Domain{{{HF}}} = \emptyset \cup \Pow{\emptyset} \cup
  \Pow{\Pow{\emptyset}} \cup \Pow{\Pow{\Pow{\emptyset}}} \cup \dots$؛
\item $\Assign{\Obj \in}{{{HF}}} = \Setabs{\tuple{x, y}}{x, y \in
  \Domain{{{HF}}}, x \in y}$۔
\end{enumerate}
\end{ex}

\begin{digress}
دا چې د څه شي د !!a{structure} ګڼلو دپاره کوم شرطونه ږدو، زموږ پر
منطق اغېز کوي۔ مثلاً د تشو ساحو نۀ منل، د سورو لرونکو جملو د صدق د
معمول حساب له مخې، دا تضمينوي چې
$\lexists[x][(!A(x) \lor \lnot !A(x))]$ معتبر دے---يعنې يو منطقي صدق
دے۔ او دا شرط چې ټول !!{constant}s بايد د ساحې کوم څيز ته اشاره وکړي،
تضمينوي چې وجودي تعميم د استنتاج يوه سمه بڼه ده:
$!A(a)$، نو $\lexists[x][!A(x)]$۔ که مو نومونو ته د ساحې بهر يا هېڅ
څيز ته اشاره کول روا ګڼلي واے، د \emph{ازاد منطق} پر لور به تللي وو؛
په هغه کښې وجودي تعميم يوه اضافي مقدمې ته اړتيا لري:
$!A(a)$ او $\lexists[x][\eq[x][a]]$، نو $\lexists[x][!A(x)]$۔
\end{digress}

\end{document}
